\section{模型假设与数据准备}
\subsection{共同假设与适用范围}
本文采用以下共同约定。第一，样本编号及附件中的对齐位置用于关联模态，预计算位置不等同于等长的物理时间；媒体回查另使用原始时间戳。第二，在题目给定的特征表示下，音视频整行全零视为无有效观测，单个坐标为零仍是合法数值，缺失状态由标准化前的掩码确定。第三，保留附件二提供的原始标签与数据划分，标准化统计量只由当前训练子集估计。问题一的时间匹配规则和问题三的交互近似在各自模型中进一步说明。
\subsection{数据来源、划分与预处理约定}
全部任务使用题目附件。附件一包含100条原始视频及文本、情感标签，样本由视频标识和片段标识联合确定。附件二提供情感强度标签与预计算特征；问题二、三统一采用aligned\_50版本，附件三、四分别用于两问的无标签专项推理。下文的训练集、验证集和测试集均指附件二给定的对应子集；训练集内部的三折划分仅用于模型选择。

\Needspace{10\baselineskip}
\begin{longtable}{lrrrrl}
\caption{数据划分与用途}\label{tab:data}\\
\toprule 集合 & 样本数 & 负向 & 中性 & 正向 & 用途\\\midrule\endfirsthead
\toprule 集合 & 样本数 & 负向 & 中性 & 正向 & 用途\\\midrule\endhead
附件二训练 &3395&967&758&1670&参数拟合\\
附件二验证 &728&206&184&338&模型比较与误差分析\\
附件二测试 &727&207&158&362&描述性评价\\
附件三 &30&—&—&—&问题二专项预测\\
附件四 &20&—&—&—&问题三预测与解释\\\bottomrule
\end{longtable}
每条预计算样本最多50个对齐位置，声学维度74、视觉维度35；文本包含768维连续表示及三行BERT输入记录。两个预测模型均使用词元编号进入BERT，训练、验证与专项推理保持接口一致。对齐位置表示附件的共同序列索引，精确物理时间通过原始媒体另行记录。

标签$y\in[-3,3]$，负值、零值、正值分别对应负向、中性、正向，类别编码$c\in\{0,1,2\}$。原始标签和划分保持不变，均值、标准差及类别频率只由相应训练子集估计。

\subsection{符号说明}
\begin{longtable}{ll}
\toprule 符号 & 含义\\\midrule
$i,j,m$ & 样本、序列位置与模态索引，$m\in\{T,A,V\}$\\
$X^m$ & 模态$m$的输入特征\\
$v,o^m,d^m$ & 有效位置、原始观测与人工缺失标记\\
$\widetilde o^m$ & 施加人工缺失后的可用位置标记\\
$y,c$ & 连续情感强度与离散极性类别\\
$\widehat y,\widehat c$ & 模型预测的强度与类别\\\bottomrule
\end{longtable}
模型可用掩码为$\widetilde o^m=v o^m(1-d^m)$。问题一另保存人脸事件F及逐维有效性；单个合法数值为零时仍保留其观测状态。

\subsection{评价指标与数据使用范围}
预测评价使用Accuracy、三类等权Macro-F1、强度MAE和Pearson相关系数：
\begin{equation}
\mathrm{Acc}=\frac1N\sum_i\mathbf1(\widehat c_i=c_i),\qquad
\mathrm{MacroF1}=\frac13\sum_{c=0}^2\frac{2TP_c}{2TP_c+FP_c+FN_c}.
\end{equation}
\begin{equation}
\mathrm{MAE}=\frac1N\sum_i|\widehat y_i-y_i|,\qquad
r=\frac{\sum_i(y_i-\bar y)(\widehat y_i-\overline{\widehat y})}{\sqrt{\sum_i(y_i-\bar y)^2\sum_i(\widehat y_i-\overline{\widehat y})^2}}.
\end{equation}
F1分母为零时该类别记0；相关系数分母为零时记未定义。Macro-F1与Accuracy评价离散类别，MAE与Pearson评价连续强度。时间候选覆盖率以原词总数为分母，解释分解误差比较贡献之和与当前模型读出，分别在相应章节报告。


问题二在训练集内进行视频分组三折选模，并在附件二验证集上比较三个随机种子的实验结果；问题三在验证集选择检查点和决策策略。历史开发中曾查看附件二测试表现，问题三最终方案也参考了已有性能与解释要求，因此测试指标用于描述当前方案，不能作为完全独立的确认性结论。数据核查还发现附件三有4条、附件四有5条与附件二测试集重合；专项集按题目要求输出预测，不回填这些样本的标签，也不重复累计为新增独立测试证据。
